\scanpage{129}
\mspage{116}
i.e. aux deux relations
\[
 h_0\widetilde i=\widetilde h_0,\qquad p_0h_0=\beta_0.\tag*{(6.3.15)}
\]
La deuxième de ces relations n'est autre que
la première relation (6.3.7), et il reste à
voir que \deleted{\unsure{6.3.}} la première relation (6.3.15)
équivaut à l'ensemble des deux relations
restantes de (6.3.7), (6.3.8) :
\[
 h_0i_0=\alpha_0,\qquad \lambda_{\deleted{\unsure{I|Y}}X}h'_{I'}=h_0\underset{\widetilde i s}{\underline{\lambda_B}}.\tag*{(6.3.16)}
\]
Or la relation en question $h_0\widetilde i=\widetilde h_0$ équivaut
à l'\unsure{ens} des deux relations
\[
\begin{aligned}
&\underbrace{\begin{array}[t]{c@{\,}c@{\,}c}
 \underline{(h_0\widetilde i)r}&=&\underline{\widetilde h_0 r},\\
 h_0(\underset{i_0}{\underline{\widetilde i r}})&&\alpha_0\\
 \underline{\phantom{h_0(\widetilde i r)}}&&\\[-2pt]
 h_0i_0&&
\end{array}}_{\substack{\text{équivalent à}\\ h_0i_0=\alpha_0}}
\quad
\underbrace{\begin{array}[t]{c@{\,}c@{\,}c}
 \underline{(\overset{\deleted{\unsure{\sim}}}{h}_0\widetilde i)s}&=&\underline{\widetilde h_0 s}\\
 \Vert&&\lambda_Xh'_{I'}\\
 h_0(\underset{\lambda_B}{\underline{\widetilde i s}})&&\\
 \underline{\phantom{h_0(\widetilde i s)}}&&\\[-2pt]
 h_0\lambda_B&&
\end{array}}_{\substack{\text{équivalent à}\\ h_0\lambda_B=\lambda_X \unsure{h''_I|h'_{I'}}}},
\end{aligned}
\]
c'est OK.

Ceci posé, dans le carré (6.3.14), les
hypothèses $i\in\Phi(I)$, $p\in\Psi\langle I\rangle$ impliquent
$\widetilde i_{\deleted{\unsure{0|2}}}\in\Phi$, $p_0\in\Psi$. Donc la flèche cherchée
$h_0$ existe, qed.

Cela achève la démonstration du cor. 6.3,
et par là \unsure{même|aussi}, celle de la dernière
partie, d), de 6.1. On est heureux !
% archive sheet number: 128

